\documentclass{article}
\usepackage{graphicx} % Required for inserting images
\usepackage{xcolor}

\usepackage{amsmath}
\usepackage{amsthm}
\usepackage{thmtools} 
\usepackage{cleveref}


% Number equations within sections (instead of chapters)
\numberwithin{equation}{section}

% Declare theorem environments anchored to sections
\declaretheorem[name=Theorem,numberwithin=section]{theorem}
\declaretheorem[name=Lemma,numberlike=theorem]{lemma}
\declaretheorem[name=Proposition,numberlike=theorem]{proposition}
\declaretheorem[name=Corollary,numberlike=theorem]{corollary}
\declaretheorem[name=Definition,numberlike=theorem]{definition}
\declaretheorem[name=Assumption,numberlike=theorem]{assumption}
\declaretheorem[name=Example,numberlike=theorem]{example}
\declaretheorem[name=Remark,numberlike=theorem]{remark}

% For restatable theorems 
\declaretheorem[name=Theorem,numberwithin=section]{restatabletheorem}

\title{Recursive Joint Simulation in Games}
\author{Colomban Duclaux}
\date{\today}

\usepackage[backend=biber,style=alphabetic]{biblatex} % Ensure style matches your preference
\addbibresource{../references.bib} % Link your .bib file

\newcommand{\dupoc}{\texttt{DUPOC}}
\newcommand{\cupod}{\texttt{CUPOD}}
\newcommand{\cimcic}{\texttt{CIMCIC}}
\newcommand{\dimcid}{\texttt{DIMCID}}
\begin{document}

\maketitle
These notes recap the content exposed in \cite{kovarik2024recursivejointsimulationgames}

\section{Abstract}
We place ourselves in the context of open-source game theory. 
The aim is to explore ways in which we can achieve cooperation.
Recursive simulation: first, agents jointly observe a simulation that is recursive before taking an action.

\section{Introduction}
Compared to humans, AI agents by their nature allow different game theoretic setups such as open-source game theory.
In this paper, we investigate the implications of the potential ability for AI agents to simulate each other during strategic interaction.
Setting: We create a chain of nested simulations of depth k, where k is determined by the probability of not allowing further simulations.
The real agents now suspect they are also in a simulation $\to$ tradeoff between seeking a direct payoff (if in reality) and adopting an indirect strategy (if in simulation)
One strategy: \textbf{grim trigger strategy}

\section{Related Work}
\begin{itemize}
    \item Infinite recursion $\to$ program equilibria \cite{critch2022cooperativeuncooperativeinstitutiondesigns}: Difference to our work: players here receive behavioral information and not a source code. Also, intermediate simulated players are also counted as rational players, whereas they would not be considered rational players in \cite{critch2022cooperativeuncooperativeinstitutiondesigns}.
    \item On program's behavior: $\epsilon$-grounded \texttt{FairBot}: separate simulation of one another instead of joint simulation.
    \item Unilateral simulation: Only one agent can simulate the other.
    \item Super rationality reasoning: Knowing that the opponent is like me so if I cooperate, they will too.
\end{itemize}

\section{Background}
\subsection{Normal-form games}
Set the foundations and notations.

\subsection{Repeated Games}
Key types of Normal form games $\mathcal{G} = (\mathcal{A}, u^\mathcal{G})$: \begin{itemize}
    \item Rep$_T(\mathcal{G})$: Players play $\mathcal{G}$ $T$ times.
    \item Rep$_\omega(\mathcal{G}, p)$: Players play infinitely many times, with exponential discounting.
    \item Rep$_{T=?}(\mathcal{G})$: players plays an unknown number of repetitions, where at each stage, the game ends with a probability of $1-p$.
    \item Rep$_T(\mathcal{G})$
\end{itemize}
Notion of \textbf{strategically equivalent} and \textbf{realization-equivalent} games.

Proposition: in an infinitely repeated Prisoner's Dilemma with sufficiently high $p$ has as an equilibrium the policy that players cooperate every single time.

Not the case in the finite setting.

\section{Recursive Joint Simulation Games and Equivalence to Repeated Games}
\begin{definition}
    A \textbf{recursive joint-simulation game} RJS$(\mathcal{G},p)$ is a normal form game $\mathcal{G}$ such that it has the recursive simulation as in the intro, with probability $p$ of ending the simulation.
\end{definition}
\begin{theorem}
    For any $\mathcal{G}$ and $p$, RJS$(\mathcal{G},p)$ is strategically equivalent to Rep$_{T=?}(\mathcal{G})$ and Rep$_\omega(\mathcal{G}, p)$.
\end{theorem}

\section{Extensions of the Formalism}
\subsection{Voluntary Simulation}
Previous section assumes decision to simulate comes from a third party. The same is the case when the players decide whether to simulate or not.
Three cases are available to extend the RJS$(\mathcal{G},p)$ setting in this case, but it becomes equivalent.

\subsection{Variable Simulation Probability}
Probability of continuing the simulation depends on the depth of the simulation. But this is realization-equivalent to Rep$_\omega(\mathcal{G}, (p_t)_t)$

\section{Self-locating Beliefs}
This section takes the point of view of the agent acting.

\section{Two Objections to the Simulation Framework}
\textbf{Indistinguishability:} Agents cannot distinguish between being in a simulation and being in the real world.
\begin{itemize}
    \item Could indistinguishability be impossible \textit{in principle}? $\to$ No because its just numerical bits for AI agents, in a simulation or the real world.
    \item Could indistinguishability be impossible \textit{in practice}? $\to$ Depending on the complexity of the simulation, we can or cannot do it in practice.
\end{itemize}

\section{Conclusion}
In some settings, recursive joint-simulation games inherit the properties of repeated games and allow a larger range of outcomes, such as cooperation.

More broadly, the setting of our paper is just one example illustrating the interplay between: (1) topics in formal epistemology, including the study of decision scenarios that in the context of ordinary human affairs might well be considered fanciful; and
(2) questions about the design of AI agents, whose strategic interactions may take a
form unlike any in ordinary human affairs.

\printbibliography
\end{document}